% Insert as A.75.  The full completed Alexander-module obstruction.
\begin{statement}{A.75}{Proved}{The full pro-\(\ell\) metabelian Fox obstruction and its finite-weight detection}
Fix a prime $\ell$, and let $F$ be the free pro-$\ell$ geometric
fundamental group of $U=\mathbf P^1_\Q\setminus\{0,1,\infty\}$,
with generators $x,y$ and rational tangent $\overrightarrow{01}$.
Put $M=F/F''$ and $D=F'/F''$, where all commutator subgroups
are closed.
For a fixed pair of continuous Kummer cocycles
$a,b:G_\Q\to\Z_\ell(1)$, the obstruction to an actual section of
$M\rtimes G_\Q$ with these abelian coordinates is the explicit
double connecting class below.
It takes values in one completed rank-one Alexander module,
with its full twisted arithmetic action.
The statement includes the exact change to the point-$3$
splitting and the local cohomology quotient at $q\ne\ell$.
The continuous map from this obstruction group to the
inverse limit of its finite-weight obstruction groups is
an isomorphism: the intervening $I$-adic $\varprojlim^1$
term vanishes.  This detects lifting for one fixed pair,
without inferring its vanishing from the depth-dependent
pairs of A.65.
\end{statement}
\begin{proof}
\emph{The complete coefficient category and the geometric map.}
Write
\[
 \Lambda=\Z_\ell[[S,T]],\quad X=1+S,\quad Y=1+T,\quad
 I=(S,T),\quad \partial=(S,T),\quad
 e=\binom{-T}{S}.
 \tag{A.75.1}
\]
The topology is the $(\ell,S,T)$-adic topology.  All modules below
are compact $\Z_\ell$ modules with jointly continuous $G_\Q$
action; cohomology means continuous inhomogeneous cochains in
this category.
The pro-$\ell$ Fox derivative, followed by abelianization, is
the continuous map
$w:F\to\Lambda^2$ characterized by
\[
 w(x)=\binom10,\quad w(y)=\binom01,\quad
 w(vv')=w(v)+\overline v\,w(v'),\qquad
                   \partial w(v)=\overline v-1 .
 \tag{A.75.2}
\]
Here $\overline v=X^aY^b$ is the group-like element recording
the image of $v$ in $F^{\rm ab}=\Z_\ell(1)^2$.
The completed Crowell sequence and its rank-two
Blanchfield--Lyndon identification are
\[
 0\longrightarrow D\xrightarrow{\,w\,}\Lambda^2
       \xrightarrow{\partial}I\longrightarrow0,\qquad
                    w(D)=\Lambda e .
 \tag{A.75.3}
\]
For these precise continuous maps see
\cite[(5.3.5)--(5.3.6), Proposition 5.3.7, and Lemma 7.1.1]
{BIR_KodaniMorishitaTerashima2017}.
The second equality can also be checked directly from the first:
$Sp+Tq=0$ implies $(p,q)=h(-T,S)$, because $S,T$ are
relatively prime in $\Lambda$.
Under our commutator convention
$[x,y]=xyx^{-1}y^{-1}$, its image is exactly $e$.

These maps identify the actual compact group $M$ with
\[
 \left\{(u,v):
    u=X^aY^b,\ v\in\Lambda^2,\ \partial v=u-1\right\},
 \qquad
 (u,v)(u',v')=(uu',v+uv').
 \tag{A.75.4}
\]
Indeed the kernel of $(\overline v,w(v))$ on $F$ is $F''$
by (A.75.3).  Conversely, lift $u$ by $x^ay^b$.
Every other vector with the same boundary differs from its Fox
vector by an element of $\Lambda e$, hence is obtained by
multiplication by an element of $D$.
This proves surjectivity; compactness proves that the resulting
continuous bijection is a homeomorphism.

\emph{The actual tangential Galois action.}
Let $\chi=\chi_\ell$ and let $\theta_g$ be the automorphism of
$\Lambda$ sending $X,Y$ to $X^{\chi(g)},Y^{\chi(g)}$.
The normalized tangential action used in A.38 is
$\rho_g(x)=x^{\chi(g)}$,
$\rho_g(y)=f_g^{-1}y^{\chi(g)}f_g$, with $f_g\in F'$;
the precise base point and this formula are those of
\cite[\S2, (2.1) and its following paragraph]
{BIR_NakamuraSchneps2000}.
Write $w(f_g)=k_g e$.
For $z\in\Z_\ell$ define the integral series
$L_X(z)=(X^z-1)/S$, $L_Y(z)=(Y^z-1)/T$.
The Fox matrix, whose columns are $w(\rho_g(x))$ and
$w(\rho_g(y))$, is
\[
 \begin{split}
 J_g&=
 \begin{pmatrix}
 L_X(\chi(g))&-T\bigl(Y^{\chi(g)}-1\bigr)k_g\\
 0&L_Y(\chi(g))+S\bigl(Y^{\chi(g)}-1\bigr)k_g
 \end{pmatrix},\\
 J_{gh}&=J_g\,\theta_g(J_h),\qquad
              \partial J_g=\theta_g(\partial),\\
 J_g\theta_g(e)&=F_g e,\qquad
 F_g=L_X(\chi(g))L_Y(\chi(g))
       +(X^{\chi(g)}-1)(Y^{\chi(g)}-1)k_g .
 \end{split}
 \tag{A.75.5}
\]
All entries are integral continuous series.
The first formula follows by differentiating
$f_g^{-1}y^{\chi(g)}f_g$ using (A.75.2).
The next two identities are the Fox chain rule and fundamental
identity.  Multiplication of the displayed matrix by
$(-\theta_g(T),\theta_g(S))^t$ gives the last identity directly.
In particular $F_g\in\Lambda^\times$ has augmentation
$\chi(g)^2$ and satisfies
$F_{gh}=F_g\theta_g(F_h)$.
This is the rank-one Ihara power series in the specified
tangential framing; the general reduced Gassner comparison is
\cite[Theorem 7.1.3]{BIR_KodaniMorishitaTerashima2017}.
We use the explicit matrix (A.75.5), so no change of base point
in that reference is left implicit.

\emph{Twisting by the fixed abelian cocycle.}
Set $u_g=X^{a(g)}Y^{b(g)}$; then
$u_{gh}=u_g\theta_g(u_h)$.
Give $\Lambda_s$, $I_s$, $P_s=\Lambda^2$, and
$D_s=\Lambda e$ the following actions:
\[
 \begin{aligned}
 g\cdot r&=u_g\theta_g(r)
                       &&(r\in\Lambda_s\text{ or }I_s),\\
 g\cdot v&=u_gJ_g\theta_g(v)
                       &&(v\in P_s),\\
 g\cdot(h e)&=u_gF_g\theta_g(h)e
                       &&(he\in D_s).
 \end{aligned}
 \tag{A.75.6}
\]
Their cocycle identities follow from (A.75.5), and the maps in
the two sequences
\[
 \begin{gathered}
 0\longrightarrow I_s\longrightarrow\Lambda_s
       \xrightarrow{\epsilon}\Z_\ell\longrightarrow0,\\
 0\longrightarrow D_s\longrightarrow P_s
       \xrightarrow{\partial}I_s\longrightarrow0
 \end{gathered}
 \tag{A.75.7}
\]
are equivariant.  The last $\Z_\ell$ has trivial action.
These sequences are split as continuous $\Z_\ell$ modules,
without requiring equivariant splittings.
For $\partial$, assign every monomial $S^iT^j$ with $i>0$
to $(S^{i-1}T^j,0)$, and those with $i=0,j>0$ to
$(0,T^{j-1})$.
Thus they really give connecting homomorphisms on continuous
cochains.

The first connecting image of $1$ is
$\eta(g)=u_g-1\in I_s$.
An explicit continuous lift is
\[
 v_g=\binom{L_X(a(g))}{X^{a(g)}L_Y(b(g))},\qquad
 \omega(g,h)e
       =v_g+u_gJ_g\theta_g(v_h)-v_{gh}.
 \tag{A.75.8}
\]
The right side has zero boundary, so the scalar $\omega(g,h)$
exists uniquely in $\Lambda$ and is continuous by (A.75.3).
It is a two-cocycle with the $D_s$ action, because it is the
cochain differential of $v$ in $P_s$.
Its class is precisely
\[
 o_\ell(a,b)=
 \delta_{(A.75.7),2}\,
 \delta_{(A.75.7),1}(1)
           =[\omega]\in H^2_{\rm cts}(G_\Q,D_s).
 \tag{A.75.9}
\]
In (A.75.4), the lifts $(u_g,v_g;g)$ have exactly the defect
$(1,\omega(g,h)e)$.
Consequently $o_\ell(a,b)=0$ is equivalent to the existence of
a continuous $h:G_\Q\to\Lambda$ with
$\delta_{D_s}h=-\omega$; the actual lifted section is then
$g\mapsto(u_g,v_g+h(g)e;g)$.
This proves both directions of the claimed criterion at the
full completed level.

\emph{Exact comparison with the point-$3$ framing.}
Choose a geometric path defining the point-$3$ section and write
$s_3(g)=z_3(g)s_{\overrightarrow{01}}(g)$.
Put $u_{3,g}=\overline{z_3(g)}$ and $w_{3,g}=w(z_3(g))$.
Its two abelian cocycles represent $\kappa_\ell(3)$ and
$\kappa_\ell(-2)$.
The point-$3$ matrices are
\[
 J_{3,g}=\bigl(u_{3,g}I_2-w_{3,g}\partial\bigr)J_g,
                    \qquad F_{3,g}=u_{3,g}F_g .
 \tag{A.75.10}
\]
Indeed differentiating $zvz^{-1}$ gives
$w(zvz^{-1})=\overline z\,w(v)+(1-\overline v)w(z)$.
This proves (A.75.10) and its original arithmetic action.
For a relative cocycle $u_{{\rm rel},g}$ the absolute cocycle is
$u_g=u_{{\rm rel},g}u_{3,g}$.
The action on $D$ is literally the same:
$u_{\rm rel}F_3\theta=uF\theta$.
If a relative lift has vector $v_{{\rm rel},g}$, its absolute
vector is
\[
            v_{{\rm abs},g}
                 =v_{{\rm rel},g}+u_{{\rm rel},g}w_{3,g}.
 \tag{A.75.11}
\]
Substitution in (A.75.8), using the cocycle identity for $z_3$,
shows that the two defect vectors are equal.
Thus this is an exact gauge comparison of cochains, not just
an identification of associated weights.

\emph{The full local coefficient module and residue quotient.}
Fix $q\ne\ell$.  The tangential action on $F$ is unramified at
$q$: the tripod and its marked tangent have good reduction
over $\Z_q$.
Discarding the prime-to-$\ell$ inertia kernel leaves
$\langle\tau\rangle\rtimes\langle\Phi\rangle
 =\Z_\ell(1)\rtimes\widehat\Z$, with
$\Phi\tau\Phi^{-1}=\tau^q$.
The discarded kernel acts trivially on (A.75.6).
Write
\[
 m=a(\tau),\quad n=b(\tau),\quad
 U_q=X^mY^n,\qquad
 \mathcal F_q(r)=u_\Phi F_\Phi\,\theta_\Phi(r),\qquad
                    \overline\Lambda_q=\Lambda/(U_q-1)\Lambda .
 \tag{A.75.12}
\]
Thus inertia on $D_s$ is exactly multiplication by $U_q$.
The semilinear map $\mathcal F_q$ preserves the displayed
ideal, because it transports $U_q$ to $U_q^q$ and
$(U_q^q-1)/(U_q-1)$ is a unit.
There is a canonical cohomology identification
\[
 H^2_{\rm cts}(G_{\Q_q},D_s)
   \simeq
 \overline\Lambda_q/
           (\mathcal F_q-q)\overline\Lambda_q .
 \tag{A.75.13}
\]
The denominator is the image of a continuous additive map;
it is closed, as the module is compact.
For a proof, continuous cohomology of the procyclic inertia
group gives
$H^1(\Z_\ell(1),D_s)=\overline\Lambda_q$ and no higher
cohomology.  Frobenius on this first cohomology is
$q^{-1}\mathcal F_q$.
Indeed a cocycle with value $v$ at $\tau$ is evaluated at
$\tau^{q^{-1}}$ in the conjugation formula, and its value
modulo $U_q-1$ is $q^{-1}v$.
The unramified procyclic group has cohomology only in
degrees zero and one, with $H^1(N)=N/(\Phi-1)N$.
These assertions for compact coefficients follow from the
resolutions of the finite quotients by the induced powers
of $(\ell,S,T)$; their finite $H^0,H^1$ groups satisfy the
required Mittag--Leffler property.
The Hochschild--Serre sequence therefore gives (A.75.13).

For completeness, the actual residue of a normalized tame
two-cocycle $c$ can be calculated before taking any weight
quotient.  Reduce its values in $\overline\Lambda_q$ and set
\[
 \begin{split}
 \psi_q(c)=q^{-1}\biggl(&
 c(\Phi,\tau)+c(\Phi\tau,\Phi^{-1})
       -c(\Phi,\Phi^{-1})\\
 &\hspace{20mm}
       -\sum_{j=1}^{q-1}c(\tau^j,\tau)\biggr)
       \quad\text{in }\overline\Lambda_q .
 \end{split}
 \tag{A.75.14}
\]
Its image modulo $q^{-1}\mathcal F_q-1$ is the class in
(A.75.13).  Multiplying the chosen lifts of
$\Phi,\tau,\Phi^{-1}$ and comparing with the $q$th power
of the lift of $\tau$ proves the formula.
Inertia acts trivially after this coefficient reduction,
so the correction involving $c(\Phi,\Phi^{-1})$ has the
displayed coefficient one.
Direct telescoping gives
$\psi_q(\delta h)=(q^{-1}\mathcal F_q-1)\overline{h(\tau)}$.
Prime-to-$\ell$ inertia can first be removed by the
coefficient-compatible averaging and normalization of A.27.
The coefficient map to $\overline\Lambda_q$ induces the
identity on inertia $H^1$, hence the residue just computed
represents the original class in (A.75.13).

For an ordinary local pair,
$(m,n)$ has one of the patterns
$(m>0,0)$, $(0,n>0)$, $(m,m)$ with $m<0$, or $(0,0)$.
The local point section makes the image of (A.75.9) in
(A.75.13) zero.
Its abelian representative can be made equal to the chosen
$a,b$ by lifting an abelian geometric conjugator to $M$.
At every nonzero valuation, however, $U_q\ne1$ in $\Lambda$:
specializing one variable to zero and using the coefficient
of the other proves this immediately.
Thus the full coefficient module of an infinite-support
ordinary ghost has nontrivial inertia at infinitely many
primes.  This is an exact action calculation, and it explains
why a theorem for an unramified-away-$\ell$ coefficient module
cannot be applied through the identity map to $D_s$.

\emph{The continuous inverse limit and vanishing of its
$I$-adic compatibility term.}
The ideals $I^rD_s$ are stable under the full action.
The maps $D_s/I^{r+1}D_s\to D_s/I^rD_s$ have continuous
$\Z_\ell$-linear sections, obtained by their monomial bases.
Hence their continuous cochain complexes are degreewise
surjective and have inverse limit $C^\bullet_{\rm cts}(G_\Q,D_s)$.
Their Milnor sequence is
\[
 \begin{split}
 0\longrightarrow&
 \varprojlim\nolimits_r^1
 H^1_{\rm cts}(G_\Q,D_s/I^rD_s)
 \longrightarrow H^2_{\rm cts}(G_\Q,D_s)\\
 &\longrightarrow
 \varprojlim_r H^2_{\rm cts}(G_\Q,D_s/I^rD_s)
 \longrightarrow0 .
 \end{split}
 \tag{A.75.15}
\]
The associated module
$I^rD_s/I^{r+1}D_s$ is
$\Z_\ell(r+2)^{r+1}$.
The quotient $M/I^rD$ is the actual metabelian quotient
of nilpotency class $r+1$, because its lower central
subgroups satisfy $\Gamma_{r+2}M=I^rD$.
Consequently the right-hand images of (A.75.9) are exactly
the fixed-depth metabelian obstructions for this same pair.

In this particular $I$-adic system the left term of
(A.75.15) is zero.  We prove this without assuming finite
ramification for the off-diagonal action.
For every integer $k\geq2$,
$H^1_{\rm cts}(G_\Q,\Z_\ell(k))$ is a finitely generated
$\Z_\ell$ module.
To see the ramification assertion first, an integral
one-cocycle restricts at $q\ne\ell$ to a homomorphism from
tame pro-$\ell$ inertia, because the coefficient action
there is trivial; wild inertia and the other tame factors
have no homomorphism to $\Z_\ell$.
Writing $t$ for its value on $\tau$, the exact tame
relation gives
$(q^k-q)t=0$, hence $t=0$.
The cocycle therefore kills every such inertia group
and its normal closure, and factors through
$G_{\Q,\{\ell,\infty\}}$.

For this fixed set of places the cohomology groups
$H^1(G_{\Q,\{\ell,\infty\}},\Z/\ell^n(k))$ are finite:
a cocycle cuts out an extension of bounded degree
unramified outside this fixed set, and Hermite finiteness
applies.  The finite $H^0$ systems have vanishing
$\varprojlim^1$, so integral $H^1$ is the inverse limit
of these finite $H^1$ groups and is compact.
The coefficient sequence for multiplication by $\ell$
injects its quotient modulo $\ell$ into the finite group
$H^1(G_{\Q,\{\ell,\infty\}},\mathbf F_\ell(k))$.
The compact form of Nakayama's lemma now proves finite
generation over $\Z_\ell$.

Each $D_s/I^rD_s$ has a finite filtration with the
Tate graded modules just displayed, all of weight at
least two.  Their underlying finite free $\Z_\ell$
coefficient sequences split continuously, so the
continuous cohomology long exact sequences apply.
Induction and the Noetherian property of $\Z_\ell$
show that every
$H^1_{\rm cts}(G_\Q,D_s/I^rD_s)$ is finitely generated
over $\Z_\ell$.
Equip these groups with their compact $\ell$-adic
module topologies; all transition maps are continuous,
being $\Z_\ell$ linear.

For any countable inverse system $(C_r,f_r)$ of compact
abelian groups with continuous maps, the algebraic
$\varprojlim^1 C_r$ is zero.
Indeed, for any $(a_r)$ the equations
$x_r-f_r(x_{r+1})=a_r$ are closed conditions in
$\prod_r C_r$.
Every finite collection is solvable by choosing the
last variable and solving backwards.
Compactness gives a simultaneous solution, which is
exactly surjectivity of the defining difference map
for $\varprojlim^1$.
Applying this argument to the preceding $H^1$ groups
proves the asserted isomorphism
\[
 H^2_{\rm cts}(G_\Q,D_s)
       \ \xrightarrow{\ \sim\ }\
 \varprojlim_r
 H^2_{\rm cts}(G_\Q,D_s/I^rD_s).
 \tag{A.75.16}
\]
Thus (A.75.9) vanishes exactly when all of its
finite-weight images for this one fixed pair vanish.
The map (A.75.16) does not remove possible
coefficient-$\ell$-power compatibility terms inside an
individual group on its right; these are different
maps, already retained in the continuous coefficient
computations A.27 and A.46.
A.71 identifies existence of a full metabelian lift of a
fixed pair with existence of full pro-$\ell$ metabelian
lifts of that same pair for every $\ell$.
No vanishing of (A.75.9) for an ordinary ghost has been
established by the comparison alone.
The exact symbolic certificate
\path{code/alexander_fox_certificate.py} and its saved output
check 20 Fox-matrix and scalar-multiplier identities, 20
semilinear composition identities, 75 inner-action gauge
identities, and 36 full scalar-defect extractions.
The completed cohomology assertions are proved above.
\end{proof}
